\part{Securing Applications and Data}

The last three units moved from a person, to a building, to a network, to a device. Unit 5 makes
the final move: into the code running on that device and the data it handles. Chapter~\ref{ch:5.1}
opens with how vulnerabilities get written into applications in the first place -- injection
flaws, buffer overflows, unvalidated input -- the software-level equivalent of the unpatched
configurations from Chapter~\ref{ch:4.1}. Chapter~\ref{ch:5.2} then applies the access-control
models first introduced in Chapter~\ref{ch:2.3} -- DAC, MAC, and RBAC -- to software permissions
and explains how a compromised low-privilege process becomes a full breach through lateral
movement when those controls are missing or misconfigured.

Chapters~\ref{ch:5.3} and \ref{ch:5.4} cover the two branches of cryptography that make the rest
of this unit's defenses possible: symmetric hashing and key storage on one side, asymmetric
(public-key) cryptography on the other. These chapters explain not just how the math works but
why an application reaches for one over the other -- hashing a password versus encrypting a
message versus verifying a certificate are three different problems with three different correct
tools. Chapter~\ref{ch:5.5} then applies all of it to the concrete case of a web application,
covering secure design practices like input validation, session management, and HTTPS that tie
directly back to the login-page trust problem opened in Chapter~\ref{ch:1.2}.

Chapter~\ref{ch:5.6} ends both the unit and the book with application threat auditing --
the software-layer version of the detection chapters that closed every previous unit. It is a
deliberate bookend: the course opened by asking how a single careless click compromises a person,
and it closes by asking how an organization audits its software to catch the compromise that
followed.
